\section{Syntax}
Expressions, types, and modes are defined by the grammar below (\(x\) ranges over variables, \(f\) over floating-point literals):
\begin{align*}
  \Omega &::= \E \mid \G \\
  \tau &::= \Unit \mid \Bool \mid \Nat \mid \Float{m} \mid \List{\tau} \mid \tau_1 \times \tau_2 \mid \tau_1 + \tau_2 \mid \tau_1 \to \tau_2 \\
  e &::= x \mid \lambda x.\,e \mid \mathsf{rec}\;f\;x.\,e \mid e_1\,e_2 \mid \langle e_1, e_2\rangle \mid \fstkw\,e \mid \sndkw\,e \mid \mathsf{inl}\,e \mid \mathsf{inr}\,e \mid \matchkw\;e\;\withkw\;\mathsf{inl}\,x \Rightarrow e_1 \mid \mathsf{inr}\,y \Rightarrow e_2 \\
    &\quad\mid [\,] \mid e_1 :: e_2 \mid \matchkw\;e\;\withkw\;[\,] \Rightarrow e_1 \mid x::xs \Rightarrow e_2 \\
    &\quad\mid \mathsf{true} \mid \mathsf{false} \mid \ifkw\;e\;\thenkw\;e_1\;\elsekw\;e_2 \mid \letkw\;x=e_1\;\inkw\;e_2 \\
    &\quad\mid c \mid -e \mid e_1 + e_2 \mid e_1 \times e_2 \mid e_1 < e_2 \mid \uniform(e_1,e_2) \mid \gaussian(e_1,e_2)
\end{align*}
We include a closed interval sampler that draws a float uniformly between the two bounds.

